\subsection{乘法群 C*}

我们从第三章 § 6 第 7 号知道，在非零复数的乘法群 \(\mathbf{C}^*\) 上诱导的拓扑与 \(\mathbf{C}^*\) 的群结构相容。由于 \(\mathbf{C}^*\) 在 \(\mathbf{C}\) 中是开的，由此得出 \(\mathbf{C}^*\) 是一个局部紧拓扑群（第一章，§ 9，第 7 号，命题 13），因而（关于乘法一致结构）是完备的；参看第三章，§ 6，第 8 号，命题 8）。乘法群 \(\mathbf{R}_+^*\) （\(>0\) 实数的）是 \(\mathbf{C}^*\) 的一个闭子群。另一个子群是绝对值为 1 的复数所成的集 \(\mathbf{U}\) ，它与单位圆周 \(\mathbf{S}_1\) （\(\mathbf{R}^2\) 的）等同，因此是一个紧群。此外：

命题 1. 拓扑群 \(\mathbf{C}^*\) 同构于拓扑群 \(\mathbf{R}^*\) 与 \(\mathbf{U}\) 的乘积。

因为映射 \(z \to \left( |z|, \frac{z}{|z|} \right)\) 是 \(\mathbf{C}^*\) 到 \(\mathbf{R}_+^* \times \mathbf{U}\) 上的一个同胚（第六章，§ 2，第 3 号，命题 3）；并且由此立即得出它是群结构的一个同构。

已经知道拓扑群 \(\mathbf{R}_{+}^{*}\) 同构于加法群 \(\mathbf{R}\) （第五章，§ 4，定理 1）；因此对拓扑群 \(\mathbf{C}^{*}\) 的研究就归结为对 \(\mathbf{U}\) 的研究，我们将在 § 2 中考虑它。

